\chapter*{Resumen}
\addcontentsline{toc}{chapter}{Resumen}

Los inversores minoristas argentinos suelen tener acciones, CEDEARs, fondos cotizados y criptoactivos repartidos en varias plataformas. Las herramientas que usan muestran el rendimiento de cada cuenta, pero no la composición de la cartera completa ni el contexto externo que la afecta, y sus alertas no explican por qué aparecen. Este Proyecto Final de Ingeniería presenta Allocia, una plataforma de análisis de portafolios multiactivo con dos clientes, web y móvil, sobre una misma cuenta y una misma API.

El sistema clasifica las posiciones con un catálogo propio, calcula pesos, exposiciones por clase, sector, región y moneda y el índice de concentración de Herfindahl--Hirschman, incorpora contexto geopolítico y señales emergentes con fuente y fecha, y genera observaciones explicables mediante un motor de reglas determinista que muestra la regla, las métricas y la evidencia de cada una. Una capa separada redacta informes por activo con un modelo de lenguaje restringido a fuentes públicas, con citas validadas literalmente y abstención cuando falta evidencia. El sistema informa y explica; no ejecuta operaciones ni recomienda comprar o vender.

La investigación con usuarios combina una encuesta a 190 personas y tres entrevistas transcriptas, que coinciden en la dispersión de la información y en la demanda de explicaciones verificables. El desarrollo sigue Scrum con un tablero Kanban, en sprints de dos semanas. La evaluación comprende 62 comprobaciones de servicios, integración, seguridad e inteligencia artificial, 33 pruebas unitarias, 27 pruebas de navegador y verificaciones sobre el servicio publicado; los cálculos coinciden con su reconstrucción manual. El Proyecto analiza además el encuadre legal (Leyes 25.326 y 26.831 y Normas de la CNV), la seguridad y la privacidad, el modelo de negocio freemium y su viabilidad en tres escenarios a 36 meses.

Los resultados cumplen el objetivo general en su núcleo: una cartera cargada desde cualquiera de los dos clientes se mide y produce observaciones cuyo fundamento puede inspeccionarse. Una evaluación con cinco participantes muestra que los inversores activos comprenden las observaciones y su evidencia, mientras que los perfiles sin experiencia requieren un lenguaje más simple; los cambios derivados se implementaron en ambos clientes.

\medskip
\noindent\textbf{Palabras clave:} análisis de portafolios, diversificación, riesgo geopolítico, inteligencia artificial explicable, motor de reglas, criptoactivos, inversor minorista.
